% !Mode:: "TeX:UTF-8"
\chapter*{结论\markboth{结论}{}}
\addcontentsline{toc}{chapter}{结论}

本文从复几何与代数两个视角梳理了 Calabi–Yau 簇的形变理论。
在光滑的紧 Kähler Calabi–Yau 流形上，利用全纯体积形式诱导的同构 \(\eta\) 和 Tian–Todorov 引理，
可以将 Maurer–Cartan 方程转化为递推的 \(\bar\partial\)‑方程，并通过 \(\partial\bar\partial\)‑引理消除所有高阶障碍，
从而证明形变空间是光滑的，即著名的 Bogomolov–Tian–Todorov 无障碍定理。

在代数方面，通过引入余切复形和 \(T^i\) 函子，我们建立了任意概形的形变障碍理论。
Schlessinger 判别法则明确了形变函子何时具有极小通用族或 pro‑可表示性。
对光滑代数簇，障碍仅由 \(H^2(\mathcal T_X)\) 给出，这与复几何的结果一致。
Ran–Kawamata 的 \(T^1\) 提升技巧进一步揭示，只要一阶形变在幂级数环的厚化下总可提升，
则整个形变空间就是形式光滑的；对于典范丛平凡的光滑射影簇，这一技巧给出了纯代数的无障碍性证明。

当簇存在奇点时，情形更为复杂。本文聚焦于只有普通双重点的三维 Calabi‑Yau 簇。
Friedman 较早地给出了用小解消的例外曲线类给出的光滑化判据：例外曲线类在 \(H^4(X,\mathbb C)\) 中的某种正线性关系是光滑化存在的充要条件，而额外要求它们生成 \(H^2(\Omega_X^2)\) 则是为了压制障碍。
后来的工作表明，上述生成条件并不必要：一旦存在一个 Kähler 解消，奇点形变就是无障碍的，
于是仅凭例外曲线类的正线性关系即可保证光滑化。这一结果深刻揭示了奇点解消的几何与形变理论的相互作用。

整体上，Calabi–Yau 簇的无障碍性是形变理论中最优美的现象之一，它与 Hodge 理论、镜像对称以及模空间的构造紧密相连。对于更一般的奇点或更高维数，形变障碍的消没机制仍然是活跃的研究课题。
